% ==========================================================================
%  Experimental Results — Continuous-Space MAPF with Flow + EPIBTShield
%  Target: ICRA 2027 (6+1 pages, IEEE two-column)
% ==========================================================================

\section{Experimental Evaluation}\label{sec:results}

We evaluate on three benchmark families of increasing difficulty.
\textbf{Set~A} (in-distribution): \texttt{random-32-32-10} and \texttt{empty-48-48}
with $N\!\in\!\{50,100\}$ agents, 25~scenarios each.
\textbf{Set~B} (near-distribution): \texttt{random-64-64-10} and
\texttt{room-32-32-4}, $N\!=\!50$.
\textbf{Set~C} (out-of-distribution): \texttt{warehouse-10-20-10-2-1},
$N\!\in\!\{50,100\}$.
All runs use a 512-step horizon unless noted; we also report 256-step
results for budget sensitivity.
The learned model is~\texttt{v4b} with 3~Euler integration steps
(ablated in Sec.~\ref{sec:ablation}).
Training seeds 42, 123, 456 are evaluated for robustness.

\paragraph{Metrics.}
\emph{AtGoal}: fraction of agents within goal tolerance at termination.
\emph{Coll}: cumulative pair-timestep collision count (lower is better).
\emph{ArrPLR}: path-length ratio restricted to agents that arrived, measuring route efficiency.
\emph{ArrStep}: mean arrival timestep of successful agents.

\paragraph{Baselines.}
ORCA~\cite{vandenberg2011orca}: velocity-obstacle planner with default parameterization.
PO-ORCA: ORCA with pairwise priority ordering to reduce symmetric deadlocks.
Straight+EPIBTShield: goal-directed constant-velocity policy passed through our shield, isolating the shield contribution.
Flow+ORCA: learned flow velocities shielded by ORCA instead of EPIBTShield, isolating the learned-policy contribution.

% ------------------------------------------------------------------
\subsection{Set~A: In-Distribution Performance}\label{sec:seta}

Table~\ref{tab:main} reports Set~A results.
On the hardest configuration (\texttt{random}, $N\!=\!100$),
Flow+EPIBTShield achieves 87.4\% AtGoal with 633~collisions,
compared with ORCA at 16.8\% AtGoal and 4169~collisions
(a 70.6 percentage-point completion gain and 84.8\% collision reduction).

The decomposition reveals that shielding is the primary mechanism:
Straight+EPIBTShield already reaches 77.4\% AtGoal,
a~60.6~pp improvement over ORCA, through priority inheritance
and velocity-space backtracking alone.
By contrast, Flow+ORCA reaches only 48.6\%---learned velocities
through a weak shield cannot recover from dense conflicts.
Adding the learned flow prior on top of EPIBTShield yields a
further 10.0~pp gain (77.4\%~$\to$~87.4\%), demonstrating that
the rectified-flow policy provides useful preferred-velocity
priors that the shield can exploit when resolving conflicts.

Averaged across all Set~A configurations at 512~steps,
Flow+EPIBTShield achieves 0.927 AtGoal at $N\!=\!50$ and 0.921 at $N\!=\!100$,
compared with 0.889/0.886 for Straight+EPIBTShield and
0.579/0.584 for ORCA (Table~\ref{tab:main}).

At 256~steps, completion drops moderately
(0.842/$N\!=\!50$, 0.813/$N\!=\!100$),
confirming that the longer horizon primarily benefits
agents in congested regions that require more timesteps to resolve.

% ------------------------------------------------------------------
\subsection{Generalization: Sets~B and~C}\label{sec:gen}

Table~\ref{tab:generalization} reports generalization results.
On near-distribution maps (Set~B), Flow+EPIBTShield maintains
its advantage: on \texttt{random-64-64-10} it reaches 61.4\% AtGoal
versus 50.5\% for Straight+EPIBTShield and 4.8\% for ORCA.
On \texttt{room-32-32-4}, a bottleneck-heavy layout, it achieves
26.2\% versus 11.4\% and 1.2\%, respectively.

On the out-of-distribution \texttt{warehouse} maps (Set~C),
a notable reversal occurs:
Straight+EPIBTShield (41.1\% at $N\!=\!50$, 39.4\% at $N\!=\!100$)
outperforms Flow+EPIBTShield (25.5\%, 21.7\%).
The learned flow prior, trained exclusively on random and empty maps,
produces velocity proposals that are poorly suited to narrow warehouse
corridors and structured shelf layouts.
EPIBTShield must repeatedly override these proposals, incurring
backtracking overhead that a simple goal-directed policy avoids.
This result underscores two points:
(1)~EPIBTShield itself generalizes robustly to unseen environments,
and (2)~the learned component requires domain-matched training data
to contribute positively.

% ------------------------------------------------------------------
\subsection{Ablation Studies}\label{sec:ablation}

\paragraph{Integration steps.}
We vary the number of Euler steps used to integrate the flow ODE
at inference (3, 5, 10, 20~steps).
As shown in Table~\ref{tab:main}, 3~steps achieves the best
completion--efficiency tradeoff.
On \texttt{random} $N\!=\!100$, AtGoal is 0.874 at 3~steps and
changes only marginally at 5~steps (0.879) before declining at
10 and 20~steps (0.856, 0.868).
Meanwhile, ArrPLR degrades monotonically: 1.462 at 3~steps
versus 1.659 at 20~steps, a~13.5\% increase in route inefficiency.
The additional integration steps produce smoother velocity fields
but also more conservative trajectories that increase path length
without reliably improving goal completion.

\paragraph{Model version.}
At 256~steps, model \texttt{v4b} consistently outperforms
\texttt{v4} across all Set~A configurations,
with gains ranging from 8.4~pp (\texttt{random} $N\!=\!100$)
to 10.6~pp (\texttt{empty} $N\!=\!50$), confirming that architectural
and training improvements in \texttt{v4b} translate to stronger
preferred-velocity proposals.

\paragraph{Multi-seed robustness.}
Table~\ref{tab:main} includes per-seed results across
seeds 42, 123, and~456.
On \texttt{empty} maps, variance is low
(mean AtGoal 0.974$\pm$0.008 at $N\!=\!50$,
0.952$\pm$0.023 at $N\!=\!100$).
On \texttt{random} maps, variance increases:
0.817$\pm$0.037 at $N\!=\!50$, 0.812$\pm$0.038 at $N\!=\!100$.
Seed~42 is the strongest across all settings;
seeds 123 and 456 still surpass Straight+EPIBTShield on
\texttt{random} $N\!=\!100$ (0.789 and 0.782 vs.\ 0.774),
but the learned margin is modest for retrained seeds.
This suggests that while the learning signal is consistent,
training stability on cluttered maps warrants further investigation
(e.g., curriculum learning or loss landscape smoothing).

% ------------------------------------------------------------------
\subsection{EPIBTShield as the Primary Contribution}\label{sec:shield-contrib}

The experimental decomposition provides clear evidence that
EPIBTShield is the primary algorithmic contribution of this work.
ORCA and PO-ORCA both fail categorically on dense random maps,
achieving under 17\% completion at $N\!=\!100$ regardless of the
velocity source.
Replacing the shield with EPIBTShield while retaining a naive
straight-line velocity policy improves completion from 16.8\%
to 77.4\%---a 4.6$\times$ increase attributable entirely to
priority inheritance and continuous-space backtracking.
The shield operates by assigning dynamic priorities to agents,
allowing higher-priority agents to commit velocities first and
requiring lower-priority agents to search a collision-free velocity
within their feasible set.
When no feasible velocity exists, the agent backtracks by yielding
its current priority, enabling deadlock resolution without global
replanning.
This mechanism transfers to unseen environments
(Set~B and warehouse maps in Set~C), where EPIBTShield delivers
the largest absolute gains over ORCA even without a learned prior.
The learned flow component provides a complementary but secondary
benefit: it biases the shield's search toward more efficient
velocity proposals, improving completion by an additional 10~pp
on in-distribution maps where training data aligns with the
evaluation environment.

% ------------------------------------------------------------------
\subsection{Contribution of the Learned Flow Policy}\label{sec:flow-contrib}

The rectified-flow GNN policy contributes by providing
spatially informed preferred velocities that encode learned
coordination patterns.
Unlike the straight-line baseline, which proposes velocities
directed rigidly toward each agent's goal, the flow policy
produces context-dependent proposals that account for nearby
agents and obstacles through the GNN's message-passing structure.
This manifests as a consistent 3--10~pp AtGoal improvement
over Straight+EPIBTShield on in-distribution and
near-distribution maps.
On \texttt{room-32-32-4}, where agents must negotiate narrow
doorways, the learned policy more than doubles the shield-only
completion rate (11.4\%~$\to$~26.2\%), suggesting that
the flow model learns useful bottleneck-avoidance behaviors.
Critically, the learned policy improves performance only when
paired with a strong shield: Flow+ORCA achieves just 48.6\%
on \texttt{random}~$N\!=\!100$, below even
Straight+EPIBTShield at 77.4\%.
The flow model does not guarantee collision avoidance; it
provides high-quality proposals that a robust conflict-resolution
mechanism can refine and enforce.

% ------------------------------------------------------------------
\subsection{Limitations}\label{sec:limitations}

Three limitations merit discussion.
First, the learned flow prior degrades on out-of-distribution
layouts: on \texttt{warehouse} maps, Flow+EPIBTShield
underperforms Straight+EPIBTShield by 15.6~pp ($N\!=\!50$)
and 17.7~pp ($N\!=\!100$).
The model, trained on random obstacle and empty grids,
produces velocity proposals misaligned with the structured
corridor geometry of warehouses, and the shield's backtracking
overhead from overriding poor proposals exceeds the cost
of computing collision-free velocities from scratch.
Incorporating warehouse-like maps in training or employing
domain adaptation techniques is a clear direction for
future work.
Second, multi-seed training variance on cluttered maps
(std of 0.038 on \texttt{random}~$N\!=\!100$) indicates
sensitivity to initialization.
While all seeds outperform Straight+EPIBTShield, the margin
is small for non-best seeds (0.8--1.5~pp), and the practical
benefit of the learned prior may be overstated if only the
best seed is reported.
We report all seeds transparently to avoid this pitfall.
Third, empirical collision counts, while substantially
reduced (84.8\% fewer than ORCA on the hardest case), remain
nonzero.
EPIBTShield does not provide a formal collision-avoidance
guarantee in continuous space under finite integration steps;
residual collisions arise from discretization of the
continuous dynamics and simultaneous velocity commitments
within a single shield iteration.
Formal safety certificates under bounded integration error
remain an open problem.


% ==================================================================
%  TABLES
% ==================================================================

\begin{table*}[t]
\centering
\caption{%
  Set~A results (in-distribution).
  \textbf{AtGoal}: fraction at goal; \textbf{Coll}: cumulative collisions;
  \textbf{ArrPLR}: path-length ratio (arrived agents);
  \textbf{ArrStep}: mean arrival step.
  Averages over 25~scenarios.  Best AtGoal per row in \textbf{bold}.
  Integration-step ablation uses Flow+EPIBTShield with $k$ Euler steps.
  Multi-seed results report per-seed AtGoal (Coll) and mean$\pm$std.
}\label{tab:main}
\vspace{2pt}
\small
\setlength{\tabcolsep}{4pt}
\begin{tabular}{@{}ll rrrr rrrr rrrr rrrr@{}}
\toprule
& & \multicolumn{4}{c}{\texttt{empty-48-48}, $N\!=\!50$}
& \multicolumn{4}{c}{\texttt{empty-48-48}, $N\!=\!100$}
& \multicolumn{4}{c}{\texttt{random-32-32-10}, $N\!=\!50$}
& \multicolumn{4}{c}{\texttt{random-32-32-10}, $N\!=\!100$} \\
\cmidrule(lr){3-6}\cmidrule(lr){7-10}\cmidrule(lr){11-14}\cmidrule(lr){15-18}
& Method & AtG & Coll & PLR & Step
       & AtG & Coll & PLR & Step
       & AtG & Coll & PLR & Step
       & AtG & Coll & PLR & Step \\
\midrule
\multirow{5}{*}{\rotatebox{90}{\scriptsize Baselines}}
& ORCA
  & .990 &   17 & 1.01 & 118
  & 1.00 &   33 & 1.01 & 147
  & .168 &  969 & 1.01 & 369
  & .168 & 4169 & 1.01 & 434 \\
& PO-ORCA
  & 1.00 &   43 & 1.05 & 125
  & 1.00 &   99 & 1.10 & 153
  & .153 &  626 & 1.10 & 413
  & .149 & 2626 & 1.10 & 443 \\
& Straight+EPIBT
  & 1.00 &    6 & 1.08 &  80
  & 1.00 &   15 & 1.10 &  97
  & .774 &  206 & 1.12 & 188
  & .774 & 1307 & 1.12 & 188 \\
& Flow+ORCA
  & ---  & ---  & ---  & ---
  & ---  & ---  & ---  & ---
  & ---  & ---  & ---  & ---
  & .486 & 3996 & 1.38 & 331 \\
\cmidrule(lr){2-18}
\multirow{2}{*}{\rotatebox{90}{\scriptsize Ours}}
& Flow+EPIBT (256)
  & .850 &   13 & 1.25 & 118
  & .812 &   80 & 1.32 & 149
  & .835 &   86 & 1.32 & 144
  & .814 &  500 & 1.38 & 158 \\
& Flow+EPIBT (512)
  & \textbf{.989} &   20 & 1.24 &  82
  & \textbf{.968} &   84 & 1.28 & 104
  & \textbf{.866} &  177 & 1.37 & 140
  & \textbf{.874} &  633 & 1.46 & 168 \\
\midrule
\multicolumn{18}{l}{\textit{Integration-step ablation} (Flow+EPIBT, 512 steps)} \\
\cmidrule(lr){2-18}
& $k\!=\!3$ (default)
  & .989 &  20 & 1.24 &  82
  & .968 &  84 & 1.28 & 104
  & .866 & 177 & 1.37 & 140
  & .874 & 633 & 1.46 & 168 \\
& $k\!=\!5$
  & .984 &  21 & 1.26 &  85
  & .961 &  96 & 1.30 & 109
  & .893 & 178 & 1.42 & 148
  & .879 & 681 & 1.59 & 175 \\
& $k\!=\!10$
  & .982 &  23 & 1.29 &  90
  & .941 & 110 & 1.34 & 115
  & .856 & 201 & 1.48 & 159
  & .856 & 744 & 1.64 & 185 \\
& $k\!=\!20$
  & .984 &  25 & 1.31 &  92
  & .956 & 105 & 1.33 & 113
  & .870 & 195 & 1.49 & 157
  & .868 & 730 & 1.66 & 183 \\
\midrule
\multicolumn{18}{l}{\textit{Multi-seed robustness} (Flow+EPIBT, 512 steps, AtGoal / Coll)} \\
\cmidrule(lr){2-18}
& Seed 42
  & .984 & --- & --- & ---
  & .981 & --- & --- & ---
  & .868 & --- & --- & ---
  & .866 & --- & --- & --- \\
& Seed 123
  & .965 & --- & --- & ---
  & .948 & --- & --- & ---
  & .781 & --- & --- & ---
  & .789 & --- & --- & --- \\
& Seed 456
  & .974 & --- & --- & ---
  & .926 & --- & --- & ---
  & .801 & --- & --- & ---
  & .782 & --- & --- & --- \\
& Mean$\pm$std
  & \multicolumn{4}{c}{.974$\pm$.008}
  & \multicolumn{4}{c}{.952$\pm$.023}
  & \multicolumn{4}{c}{.817$\pm$.037}
  & \multicolumn{4}{c}{.812$\pm$.038} \\
\bottomrule
\end{tabular}
\end{table*}


\begin{table}[t]
\centering
\caption{%
  Generalization results (Sets~B and~C).
  Set~B maps are near-distribution; \texttt{warehouse} is
  out-of-distribution.  All runs use 512~steps, 25~scenarios.
  Best AtGoal per row in \textbf{bold}.
}\label{tab:generalization}
\vspace{2pt}
\small
\setlength{\tabcolsep}{3.5pt}
\begin{tabular}{@{}l l rrr@{}}
\toprule
Map & Method & AtGoal & Coll & ArrPLR \\
\midrule
\multirow{3}{*}{\shortstack[l]{\texttt{random-64}\\$N\!=\!50$}}
  & ORCA             & .048 &  383 & 0.99 \\
  & Straight+EPIBT   & .505 &  267 & 1.05 \\
  & Flow+EPIBT       & \textbf{.614} &  106 & 1.39 \\
\midrule
\multirow{3}{*}{\shortstack[l]{\texttt{room-32}\\$N\!=\!50$}}
  & ORCA             & .012 & 1028 & --- \\
  & Straight+EPIBT   & .114 & 2020 & --- \\
  & Flow+EPIBT       & \textbf{.262} & 1113 & --- \\
\midrule
\multirow{3}{*}{\shortstack[l]{\texttt{warehouse}\\$N\!=\!50$}}
  & ORCA             & .132 &  855 & --- \\
  & Straight+EPIBT   & \textbf{.411} &   30 & --- \\
  & Flow+EPIBT       & .255 &   57 & --- \\
\midrule
\multirow{3}{*}{\shortstack[l]{\texttt{warehouse}\\$N\!=\!100$}}
  & ORCA             & .125 & 3496 & --- \\
  & Straight+EPIBT   & \textbf{.394} &  168 & --- \\
  & Flow+EPIBT       & .217 &  320 & --- \\
\bottomrule
\end{tabular}
\end{table}
